L'isòtop radioactiu fluor 18 es fa servir com a radiofàrmac en tomografies per emissió de positrons (TEP). Quan es desintegra radioactivament, aquest isòtop desprèn un positró que s'anihila ràpidament amb un electró de l'entorn i produeix dos fotons gamma amb la mateixa energia. Aquests fotons, detectats per l'aparell mèdic, permeten obtenir imatges útils per a la diagnosi. El període de semidesintegració del fluor 18 és de 109,77 minuts i podem escriure l'equació de la desintegració de la manera següent: ${}^{18}_{9}\mathrm{F} \rightarrow {}^{A}_{B}\mathrm{Y} + {}^{C}_{D}\text{positró} + {}^{0}_{0}\nu_\mathrm{e}$, en què Y és el nucli fill i $\nu_\mathrm{e}$ és un neutrí electrònic.

\figura[0.45]{fig1.png}

\begin{apartat}{1}
Indiqueu quants protons i quants neutrons té el nucli de fluor 18. Calculeu els coeficients $A$, $B$, $C$ i $D$ de l'equació i la freqüència dels fotons gamma detectats per l'aparell de la tomografia.
\end{apartat}

\begin{apartat}{1}
Calculeu el temps que ha de transcórrer perquè el nombre de nuclis de fluor 18 que queden sense desintegrar en el cos del pacient sigui l'1\,\% dels que hi havia a l'inici de la prova.
\end{apartat}

\dades{%
  Velocitat de la llum, $c = 3,00\times 10^{8}\ \mathrm{m\,s^{-1}}$.\\
  $m_\text{electró} = 9,11\times 10^{-31}\ \mathrm{kg}$.\\
  Constant de Planck, $h = 6,63\times 10^{-34}\ \mathrm{J\,s}$.}
